% !TeX program = xelatex
\documentclass[11pt, a4paper]{article}

\usepackage[utf8]{inputenc}
\usepackage{geometry}
\geometry{a4paper, left=15mm, right=15mm, top=20mm, bottom=20mm}
\usepackage{amsmath, amssymb, amsthm}
\usepackage{physics}
\usepackage{graphicx}
\usepackage{xcolor}
\usepackage{tcolorbox}
\tcbuselibrary{skins, breakable}
\usepackage{enumitem}
\usepackage{fancyhdr}
\usepackage{lastpage}
\usepackage{tabularx}
\usepackage{booktabs}
\usepackage{xeCJK}
\setCJKmainfont{SimSun}

\definecolor{navy}{RGB}{20, 30, 80}
\definecolor{slate}{RGB}{70, 80, 100}
\definecolor{royal}{RGB}{30, 80, 180}
\definecolor{forest}{RGB}{40, 120, 60}
\definecolor{crimson}{RGB}{180, 40, 40}
\definecolor{border}{RGB}{200, 205, 215}
\definecolor{paper}{RGB}{250, 251, 253}
\definecolor{mist}{RGB}{140, 150, 160}

\newcommand{\alerttag}[1]{\textsf{\textbf{\color{crimson}[#1]}}}
\newcommand{\checkbox}{$\square$}

\newtcolorbox{briefing}[1]{
  enhanced, breakable,
  colback=white, colframe=navy, boxrule=1pt, arc=3mm,
  fonttitle=\sffamily\bfseries\large, title={#1},
  colbacktitle=navy, coltitle=white,
  attach boxed title to top left={yshift=-2mm, xshift=5mm},
  boxed title style={boxrule=0pt, arc=2mm}
}

\newtcolorbox{problem}[2]{
  enhanced, breakable,
  colback=white, colframe=royal, boxrule=0.8pt, arc=2mm,
  left=4mm, right=4mm, top=3mm, bottom=3mm,
  fonttitle=\sffamily\bfseries, title={#1},
  colbacktitle=paper, coltitle=royal,
  titlerule=0pt, toptitle=2mm, bottomtitle=2mm,
  after title={\hfill\textmd{\color{slate}#2}}
}

\newtcolorbox{solution}[1]{
  enhanced, breakable,
  colback=paper, colframe=forest, boxrule=0.8pt, arc=2mm,
  fonttitle=\sffamily\bfseries, title={#1},
  colbacktitle=forest, coltitle=white,
  attach boxed title to top left={yshift=-2mm, xshift=5mm},
  boxed title style={boxrule=0pt, arc=2mm},
  top=4mm
}

\newtcolorbox{debrief}[1]{
  enhanced, breakable,
  colback=white, colframe=slate, boxrule=1pt, arc=3mm,
  fonttitle=\sffamily\bfseries\large, title={#1},
  colbacktitle=slate, coltitle=white,
  attach boxed title to top left={yshift=-2mm, xshift=5mm},
  boxed title style={boxrule=0pt, arc=2mm}
}

\pagestyle{fancy}
\fancyhf{}
\fancyhead[L]{\sffamily\color{slate} 2026 UCAS Daily Tournament}
\fancyhead[R]{\sffamily\color{slate} \thepage\ / \pageref{LastPage}}
\renewcommand{\headrulewidth}{0.4pt}
\renewcommand{\headrule}{\hbox to\headwidth{\color{border}\leaders\hrule height \headrulewidth\hfill}}

\begin{document}

\begin{center}
  {\Huge\sffamily\bfseries\color{navy} 每日大赛 \textbar\ Daily Tournament v6.0}\\[6pt]
  {\small\sffamily\color{mist} 2026-09-12 \quad\textbar\quad Phase 2: 地毯式基础回填 (Week 3-4)}\\[-3pt]
  {\color{navy}\rule{\textwidth}{0.8pt}}
\end{center}

\vspace{-4pt}

\begin{briefing}{\S 1\quad Tactical Briefing \& Warm-Up}
\begin{itemize}[leftmargin=1.5em, itemsep=2pt]
  \item \textbf{604}: 回填极限计算（泰勒公式定阶相消）与常微分方程基础。
  \item \textbf{811}: 退回物理直觉层，攻克 Heisenberg 图像的运动方程与算符演化。
\end{itemize}
\textbf{Warm-Up Retrieval} (5\,min): 在草稿纸上盲推：\\
\checkbox\ 常见函数 $e^x, \sin x, \cos x, \ln(1+x)$ 展开到第三阶 \quad 
\checkbox\ 算符 Heisenberg 运动方程公式
\end{briefing}

\vspace{10pt}

% ================= Section A =================
{\Large\sffamily\bfseries\color{navy} \S A\quad 基础肌肉记忆区 (Muscle Memory Matrix)}
\hfill{\small\itshape 604 \& 811 Foundation (30-40 min)}
\vspace{4pt}

\begin{problem}{A.1 \textbar\ Calculus: Taylor Limit}{Suggested time: 10\,min}
利用泰勒公式（带佩亚诺余项）求下列极限：
\[ \lim_{x \to 0} \frac{\sin x - x \cos x}{x^3} \]
\end{problem}
\vspace{4pt}

\begin{problem}{A.2 \textbar\ Linear Algebra: Eigenvalues}{Suggested time: 10\,min}
求实对称矩阵 $A = \begin{pmatrix} 2 & 0 & 1 \\ 0 & 2 & 0 \\ 1 & 0 & 2 \end{pmatrix}$ 的全部特征值。
\end{problem}
\vspace{4pt}

\begin{problem}{A.3 \textbar\ ODE Foundation}{Suggested time: 10\,min}
求解一阶线性非齐次常微分方程：
\[ y' - 2xy = 1 \]
（只需写出通解的积分形式，若积分不可积可保留积分号）。
\end{problem}

\clearpage
% ================= Section B =================
{\Large\sffamily\bfseries\color{navy} \S B\quad 量子概念闪击战 (Quantum Concept Blitz)}
\hfill{\small\itshape Physical Intuition without Heavy Math (10 min)}
\vspace{4pt}

\begin{problem}{B.1 \textbar\ Conceptual True/False}{Suggested time: 5\,min}
判断正误并简述理由：\\
“在 Heisenberg 图像中，系统的态矢 $|\psi\rangle$ 会随时间演化，而力学量算符 $\hat{A}$ 保持不含时。”
\end{problem}
\vspace{4pt}

\begin{problem}{B.2 \textbar\ Physical Intuition: Probability Current}{Suggested time: 5\,min}
量子力学中推导出连续性方程 $\nabla \cdot \vec{j} + \frac{\partial \rho}{\partial t} = 0$。\\
请用大白话解释这个方程的物理意义，并说明 $\vec{j}$ 是什么。
\end{problem}

\vspace{10pt}
% ================= Section C =================
{\Large\sffamily\bfseries\color{navy} \S C\quad 核心母题攻坚 (Core Model Mastery)}
\hfill{\small\itshape Standard Exam Model (20-30 min)}
\vspace{4pt}

\begin{problem}{C.1 \textbar\ Heisenberg Equation of Motion}{Suggested time: 20\,min}
考虑一维谐振子，其哈密顿量为 $\hat{H} = \frac{\hat{p}^2}{2m} + \frac{1}{2}m\omega^2 \hat{x}^2$。\\
在 Heisenberg 图像中，算符随时间的演化满足：
\[ \frac{d\hat{A}_H}{dt} = \frac{1}{i\hbar}[\hat{A}_H, \hat{H}] \]
请分别计算坐标算符 $\hat{x}_H(t)$ 和动量算符 $\hat{p}_H(t)$ 的 Heisenberg 运动方程（即求出 $d\hat{x}_H/dt$ 和 $d\hat{p}_H/dt$），并指出它们与经典牛顿力学方程的关系（Ehrenfest 定理的体现）。\\
\textit{提示：利用基础对易关系 $[\hat{x}, \hat{p}] = i\hbar$。}
\end{problem}

\clearpage
% ================= Section D =================
{\Large\sffamily\bfseries\color{navy} \S D\quad 极限挑战区 (The Abyss - 选做)}
\hfill{\small\itshape True Exam Difficulty \& Beyond (Optional, 15+ min)}
\vspace{4pt}

\begin{problem}{D.1 \textbar\ Gaussian Integral Trick}{Suggested time: 15\,min}
求解著名的高斯积分（量子力学中最常用的积分之一）：
\[ I = \int_{-\infty}^{\infty} e^{-ax^2} dx \quad (a > 0) \]
\textit{提示：考虑构造 $I^2 = \int_{-\infty}^{\infty} e^{-ax^2} dx \int_{-\infty}^{\infty} e^{-ay^2} dy$，并将其转化为二重积分的极坐标形式。这是二重积分几何定限技巧的巅峰应用。}
\end{problem}

\clearpage
% ================= Solutions =================
{\Large\sffamily\bfseries\color{navy} \S E\quad Model Solutions \& Commentary}
\vspace{4pt}

\begin{solution}{Section A \textperiodcentered\ Solutions}
\textbf{A.1} 根据泰勒公式：$\cos x = 1 - \frac{x^2}{2!} + o(x^2)$，$\sin x = x - \frac{x^3}{3!} + o(x^3)$。\\
分子：$\sin x - x\cos x = (x - \frac{x^3}{6} + o(x^3)) - x(1 - \frac{x^2}{2} + o(x^2)) = x - \frac{x^3}{6} - x + \frac{x^3}{2} + o(x^3) = \frac{x^3}{3} + o(x^3)$。\\
极限为 $\lim_{x \to 0} \frac{x^3/3}{x^3} = \frac{1}{3}$。\\
\alerttag{避坑} 遇到相减，一定要把两项展开到同次幂且非零。

\vspace{2mm}
\textbf{A.2} 写出特征多项式 $|\lambda I - A| = \begin{vmatrix} \lambda-2 & 0 & -1 \\ 0 & \lambda-2 & 0 \\ -1 & 0 & \lambda-2 \end{vmatrix}$。\\
按第二行展开：$(\lambda-2) [(\lambda-2)^2 - 1] = 0$。\\
解得 $\lambda_1 = 2$，$\lambda_2 = 3$，$\lambda_3 = 1$。

\vspace{2mm}
\textbf{A.3} 一阶线性非齐次方程 $y' + p(x)y = q(x)$。此处 $p(x) = -2x, q(x) = 1$。\\
积分因子 $I(x) = e^{\int -2x dx} = e^{-x^2}$。\\
两边同乘积分因子：$(ye^{-x^2})' = e^{-x^2}$。\\
两边积分：$ye^{-x^2} = \int e^{-x^2} dx + C$。通解为 $y = e^{x^2} (\int e^{-x^2} dx + C)$。
\end{solution}

\begin{solution}{Section B \textperiodcentered\ Solutions}
\textbf{B.1} \textbf{错误 (False)}。正好说反了。在 Heisenberg 图像中，系统的态矢固定在初始时刻 $|\psi_H\rangle = |\psi(0)\rangle$ 不随时间变化，而力学量算符随时间演化，包含了全部的动力学信息。题干描述的是 Schrödinger 图像。

\vspace{2mm}
\textbf{B.2} 这就是量子力学版本的“电荷守恒定律”。$\rho = |\psi|^2$ 是某点发现粒子的概率密度。该方程说明：空间某一点处粒子出现概率的减少速率，等于粒子以概率流 $\vec{j}$ 流出该点的散度。这保证了全空间的概率永远归一化为 1。
\end{solution}

\begin{solution}{Section C \textperiodcentered\ Solutions}
\textbf{Derivation:}\\
1. 求 $\hat{x}_H$ 的运动方程：\\
利用交换子公式 $[A, BC] = [A,B]C + B[A,C]$。\\
$[\hat{x}, \hat{H}] = [\hat{x}, \frac{\hat{p}^2}{2m} + \frac{1}{2}m\omega^2 \hat{x}^2] = \frac{1}{2m}[\hat{x}, \hat{p}^2] + 0 = \frac{1}{2m}([\hat{x}, \hat{p}]\hat{p} + \hat{p}[\hat{x}, \hat{p}]) = \frac{i\hbar}{m}\hat{p}$。\\
所以 $\frac{d\hat{x}_H}{dt} = \frac{1}{i\hbar} ( \frac{i\hbar}{m}\hat{p}_H ) = \frac{\hat{p}_H}{m}$。

2. 求 $\hat{p}_H$ 的运动方程：\\
$[\hat{p}, \hat{H}] = [\hat{p}, \frac{1}{2}m\omega^2 \hat{x}^2] = \frac{1}{2}m\omega^2 ([\hat{p}, \hat{x}]\hat{x} + \hat{x}[\hat{p}, \hat{x}]) = \frac{1}{2}m\omega^2 (-2i\hbar\hat{x}) = -i\hbar m\omega^2 \hat{x}$。\\
所以 $\frac{d\hat{p}_H}{dt} = \frac{1}{i\hbar} ( -i\hbar m\omega^2 \hat{x}_H ) = -m\omega^2 \hat{x}_H$。

\textbf{结论}: 这两个算符的演化方程形式与经典力学中谐振子的正则方程 $\dot{x} = p/m$ 和 $\dot{p} = -m\omega^2 x$ 完全一致！这就是 Ehrenfest 定理最优雅的体现。
\end{solution}

\begin{solution}{Section D \textperiodcentered\ Solutions}
\textbf{Derivation:}\\
$I^2 = \iint_{\mathbb{R}^2} e^{-a(x^2+y^2)} dx dy$。\\
转换为极坐标：$x^2+y^2 = r^2$，$dx dy = r dr d\theta$。\\
区域是整个平面，故 $r \in [0, \infty), \theta \in [0, 2\pi]$。\\
$I^2 = \int_0^{2\pi} d\theta \int_0^\infty e^{-ar^2} r dr = 2\pi [ -\frac{1}{2a}e^{-ar^2} ]_0^\infty = 2\pi (0 - (-\frac{1}{2a})) = \frac{\pi}{a}$。\\
由于被积函数为正，$I > 0$，故 $I = \sqrt{\frac{\pi}{a}}$。
\end{solution}

\clearpage
% ================= Debrief =================
{\Large\sffamily\bfseries\color{navy} \S F\quad Daily Debrief \& XP Reward}
\vspace{6pt}

\begin{debrief}{Post-Match Review (XP Checkout)}
\renewcommand{\arraystretch}{1.3}
\sffamily\small
\begin{tabularx}{\linewidth}{@{}p{3cm}X@{}}
  \textbf{Section A (肌肉记忆)} & \checkbox\ 基础全对 (+30 XP) \quad \checkbox\ 泰勒漏阶错漏 \quad \checkbox\ 积分遗忘\\
  \midrule
  \textbf{Section B (概念闪击)} & \checkbox\ 物理直觉清晰 (+10 XP) \quad \checkbox\ 图像混淆\\
  \midrule
  \textbf{Section C (核心攻坚)} & \checkbox\ 独立推导算符方程 (+50 XP) \quad \checkbox\ 依赖解析 \quad \checkbox\ 交换子计算错误\\
  \midrule
  \textbf{Section D (深渊挑战)} & \checkbox\ 成功应用极坐标双重积分 (+20 XP, 选做成就) \quad \checkbox\ 未尝试\\
  \midrule
  \textbf{Total XP Earned} & \textbf{\underline{\hspace{2cm}} XP} \\
\end{tabularx}
\end{debrief}

\vspace{10pt}

\begin{tcolorbox}[
  enhanced, colback=white, colframe=border, boxrule=0.6pt, arc=2mm,
  left=4mm, right=4mm, top=3.5mm, bottom=4mm,
  fonttitle=\sffamily\bfseries\color{slate},
  title={\raisebox{0pt}{\color{forest}\rule{3pt}{10pt}}\hspace{4pt}Key Derivation Insights \& Traps \textbar\ 突破口与避坑沉淀},
  colbacktitle=paper, coltitle=slate,
  titlerule=0pt, toptitle=2mm, bottomtitle=2mm
]
\sffamily\small
\textbf{Core Breakthrough (今日核心突破口与关键一步)}:\\[30pt]
\rule{\linewidth}{0.3pt}\\[12pt]
\textbf{Fatal Trap \& Common Error (致命陷阱与避坑警告)}:\\[30pt]
\rule{\linewidth}{0.3pt}\\[12pt]
\textbf{Day +3 Action Item (3 天后间隔重做提取的具体题号与断点)}:\\[25pt]
\rule{\linewidth}{0.3pt}
\end{tcolorbox}

\end{document}
